\section{Future Directions}

\subsection{Cryptographic Protocol}

The present paper stops at the mathematical and formalization boundary. A
protocol paper should turn the commitment, oracle, proof-carrying, and
proximity surfaces into explicit adversary games and executable verification
paths. The relevant Lean surfaces are already separated: \texttt{TCB} records
proof-carrying artifacts and runtime obligations, \texttt{Gates} records
object-local admissibility predicates, and \texttt{Consensus} records the
Proof-of-Value scoring and policy layer. The concrete protocol must specify
which certificates are public evidence, which witnesses remain private, and
which runtime failures cause rejection, hold, or re-evaluation.

\subsection{Anonymous Identity System}

The hidden-representative/public-classifier split suggests an anonymous
identity architecture. The public identity surface would expose class values,
hashes, and oracle judgments, while the private side would retain fibration
representatives, transport data, and the Hodge operator selecting the operative
kernel. In Lean, this direction is anchored by \texttt{Axioms.HomotopyClass},
\texttt{Axioms.Fibration}, and \texttt{Axioms.OracleModality}. A future system
should make the unlinkability, class-membership, and nullifier properties
explicit rather than treating public-shadow equality as strict identity.

\subsection{Sheaf Storage}

The commitment layer can also be developed as a storage and retrieval system
for controlled sheaf/cochain projections. The public store would retain
topology hashes, cochain roots, compressed shadows, and proof-carrying
statements. Full sheaf and fibration witnesses would remain private or be
opened only through dedicated proof paths. The relevant Lean interfaces are
\texttt{Axioms.CochainCommitment}, \texttt{TCB.CompressionBoundary},
\texttt{TCB.ObservationBoundary}, and \texttt{Proximity.Public}.

\subsection{Vertra Chain}

Vertra is the working name for the chain/network layer. The paper should keep
the mathematical protocol name separate from product and token naming.
The chain layer belongs in the sequel: it consumes the mathematical surfaces
from this paper as protocol evidence, but it also requires block validity,
ledger updates, fork choice, oracle attestation, worker reliability, and
runtime trust-boundary analysis. Lean files such as \texttt{Chain},
\texttt{Consensus}, and \texttt{Consensus.ProofOfValue} already mark this
boundary without making those architecture claims part of the HCP theorem.

\subsection{Complexity-Theoretic Extensions}

The current formalization isolates the hidden-operator obstruction: public data
does not canonically determine the operative Hodge--Laplacian kernel. Future
work should refine this into sharper complexity assumptions, adversary classes,
and oracle-separation statements. Natural targets include quantitative families
of hidden kernels, reductions between HCP variants, leakage models for public
classifiers, and a clearer boundary between external CVP assumptions and
project-native hidden-operator ambiguity.
